% C Shim Architecture: Existence Proof and Functional Audit
% Compile with: pdflatex C_shim_audit.tex (run twice for cross-references)
\documentclass[11pt,a4paper]{article}

% ─── Page geometry ───────────────────────────────────────────────────────────
\usepackage[
  top=25mm, bottom=25mm, left=28mm, right=28mm,
  headheight=14pt
]{geometry}

% ─── Typography ──────────────────────────────────────────────────────────────
\usepackage[T1]{fontenc}
\usepackage{lmodern}
\usepackage[scaled=0.85]{sourcecodepro}  % monospace font for code
\usepackage{microtype}                    % microtypographic refinements

% ─── Colors ──────────────────────────────────────────────────────────────────
\usepackage[dvipsnames,svgnames,x11names]{xcolor}

\definecolor{codeBackground}{HTML}{F5F5F5}
\definecolor{codeBorder}{HTML}{CCCCCC}
\definecolor{codeKeyword}{HTML}{0000CC}
\definecolor{codeComment}{HTML}{008000}
\definecolor{codeString}{HTML}{A31515}
\definecolor{codeType}{HTML}{267F99}
\definecolor{shellBackground}{HTML}{1E1E1E}
\definecolor{shellForeground}{HTML}{D4D4D4}
\definecolor{shellPrompt}{HTML}{6A9955}
\definecolor{linkColor}{HTML}{0366D6}
\definecolor{findingLow}{HTML}{FFF3CD}
\definecolor{findingMedium}{HTML}{FFE0B2}
\definecolor{findingPositive}{HTML}{D4EDDA}
\definecolor{patternColor}{HTML}{E3F2FD}
\definecolor{quoteBackground}{HTML}{F0F4F8}
\definecolor{quoteBorder}{HTML}{4A90D9}
\definecolor{passColor}{HTML}{28A745}
\definecolor{partialColor}{HTML}{FFC107}

% ─── Code listings ───────────────────────────────────────────────────────────
\usepackage{listings}

\lstdefinestyle{cstyle}{
  language=C,
  basicstyle=\ttfamily\small,
  keywordstyle=\color{codeKeyword}\bfseries,
  commentstyle=\color{codeComment}\itshape,
  stringstyle=\color{codeString},
  backgroundcolor=\color{codeBackground},
  frame=single,
  rulecolor=\color{codeBorder},
  framesep=4pt,
  xleftmargin=6pt,
  xrightmargin=6pt,
  breaklines=true,
  breakatwhitespace=false,
  showstringspaces=false,
  tabsize=4,
  columns=flexible,
  morekeywords={z_result_t,z_owned_session_t,z_moved_config_t,z_open_options_t,
    z_loaned_session_t,z_loaned_keyexpr_t,z_moved_bytes_t,z_put_options_t,
    z_moved_config_t,z_moved_session_t,z_moved_publisher_t,z_moved_subscriber_t,
    z_moved_encoding_t,z_owned_config_t,z_owned_bytes_t,z_owned_publisher_t,
    z_owned_subscriber_t,z_owned_encoding_t,z_owned_string_t,z_owned_shm_mut_t,
    z_owned_shm_provider_t,z_moved_shm_mut_t,z_moved_shm_provider_t,
    z_moved_string_t,z_moved_closure_sample_t,z_subscriber_options_t,
    z_publisher_put_options_t,z_publisher_options_t,z_delete_options_t,
    z_owned_closure_sample_t,z_loaned_sample_t,z_loaned_publisher_t,
    z_view_keyexpr_t,z_view_string_t,z_id_t,z_buf_layout_alloc_result_t,
    z_congestion_control_t,z_priority_t,z_reliability_t,z_loaned_timestamp_t,
    z_moved_source_info_t,z_scout_options_t,z_owned_closure_hello_t,
    z_owned_closure_matching_status_t,z_owned_closure_zid_t,
    ZENOHC_API,FFI_PLUGIN_EXPORT,
    Dart_CObject,Dart_Port_DL,
    size_t,int64_t,intptr_t,uint8_t,bool},
  literate={->}{$\rightarrow$}1
           {<-}{$\leftarrow$}1,
  aboveskip=\medskipamount,
  belowskip=\medskipamount,
}

\lstdefinestyle{dartstyle}{
  language=Java,  % close enough for syntax highlighting
  basicstyle=\ttfamily\small,
  keywordstyle=\color{codeKeyword}\bfseries,
  commentstyle=\color{codeComment}\itshape,
  stringstyle=\color{codeString},
  backgroundcolor=\color{codeBackground},
  frame=single,
  rulecolor=\color{codeBorder},
  framesep=4pt,
  xleftmargin=6pt,
  xrightmargin=6pt,
  breaklines=true,
  breakatwhitespace=false,
  showstringspaces=false,
  tabsize=2,
  columns=flexible,
  morekeywords={final,var,void,dynamic,String,int,bool,Pointer,Void,Uint8,
    StreamController,ReceivePort,Sample,SampleKind,Uint8List,ZBytes,Config,
    Session,Subscriber,Publisher,ZenohException,StateError,Completer},
  aboveskip=\medskipamount,
  belowskip=\medskipamount,
}

\lstdefinestyle{shellstyle}{
  basicstyle=\ttfamily\small\color{shellForeground},
  backgroundcolor=\color{shellBackground},
  frame=single,
  rulecolor=\color{shellBackground},
  framesep=4pt,
  xleftmargin=6pt,
  xrightmargin=6pt,
  breaklines=true,
  columns=flexible,
  aboveskip=\medskipamount,
  belowskip=\medskipamount,
}

\lstdefinestyle{yamlstyle}{
  basicstyle=\ttfamily\small,
  keywordstyle=\color{codeKeyword},
  commentstyle=\color{codeComment}\itshape,
  backgroundcolor=\color{codeBackground},
  frame=single,
  rulecolor=\color{codeBorder},
  framesep=4pt,
  xleftmargin=6pt,
  xrightmargin=6pt,
  breaklines=true,
  columns=flexible,
  aboveskip=\medskipamount,
  belowskip=\medskipamount,
}

% ─── Tables ──────────────────────────────────────────────────────────────────
\usepackage{booktabs}
\usepackage{tabularx}
\usepackage{longtable}
\usepackage{multirow}
\usepackage{array}

\newcolumntype{C}{>{\centering\arraybackslash}X}
\newcolumntype{L}{>{\raggedright\arraybackslash}X}
\newcolumntype{Y}[1]{>{\centering\arraybackslash}p{#1}}

% ─── Boxes and callouts ─────────────────────────────────────────────────────
\usepackage[most]{tcolorbox}

\newtcolorbox{quotebox}{
  colback=quoteBackground,
  colframe=quoteBorder,
  boxrule=0pt,
  leftrule=3pt,
  arc=0pt,
  outer arc=0pt,
  left=8pt, right=8pt, top=6pt, bottom=6pt,
  fontupper=\itshape\small,
}

\newtcolorbox{findingbox}[2][]{
  colback=#2,
  colframe=#2!60!black,
  boxrule=0.5pt,
  arc=2pt,
  left=8pt, right=8pt, top=6pt, bottom=6pt,
  fonttitle=\bfseries\small,
  title={#1},
}

\newtcolorbox{patternbox}{
  colback=patternColor,
  colframe=patternColor!50!black,
  boxrule=0.5pt,
  arc=2pt,
  left=8pt, right=8pt, top=6pt, bottom=6pt,
}

% ─── Lists ───────────────────────────────────────────────────────────────────
\usepackage{enumitem}
\setlist{noitemsep, topsep=4pt}

% ─── Section formatting ─────────────────────────────────────────────────────
\usepackage{titlesec}
\titleformat{\section}{\Large\bfseries\sffamily}{\thesection}{1em}{}
\titleformat{\subsection}{\large\bfseries\sffamily}{\thesubsection}{1em}{}
\titleformat{\subsubsection}{\normalsize\bfseries\sffamily}{\thesubsubsection}{1em}{}
\titlespacing*{\section}{0pt}{18pt plus 4pt minus 2pt}{8pt plus 2pt}
\titlespacing*{\subsection}{0pt}{14pt plus 3pt minus 2pt}{6pt plus 2pt}
\titlespacing*{\subsubsection}{0pt}{10pt plus 2pt minus 1pt}{4pt plus 1pt}

% ─── Headers and footers ────────────────────────────────────────────────────
\usepackage{fancyhdr}
\pagestyle{fancy}
\fancyhf{}
\fancyhead[L]{\small\sffamily C Shim Architecture Audit}
\fancyhead[R]{\small\sffamily zenoh-dart / zenoh-c v1.7.2}
\fancyfoot[C]{\small\thepage}
\renewcommand{\headrulewidth}{0.4pt}
\renewcommand{\footrulewidth}{0pt}

% ─── Hyperlinks ──────────────────────────────────────────────────────────────
\usepackage[
  colorlinks=true,
  linkcolor=linkColor,
  urlcolor=linkColor,
  citecolor=linkColor,
  bookmarks=true,
  bookmarksnumbered=true,
  pdfauthor={zenoh-dart Project},
  pdftitle={C Shim Architecture: Existence Proof and Functional Audit},
  pdfsubject={Dart FFI / zenoh-c v1.7.2 C Shim Layer Analysis},
]{hyperref}

% ─── Convenience commands ────────────────────────────────────────────────────
\newcommand{\code}[1]{\texttt{#1}}
\newcommand{\file}[1]{\texttt{#1}}
\newcommand{\fn}[1]{\texttt{#1}}
\newcommand{\type}[1]{\texttt{#1}}
\newcommand{\var}[1]{\texttt{#1}}
\newcommand{\PASS}{\textcolor{passColor}{\textbf{PASS}}}
\newcommand{\PARTIAL}{\textcolor{partialColor}{\textbf{PARTIAL}}}
\newcommand{\pattern}[1]{\textbf{Pattern~#1}}

% Checkmark and cross for the matrix
\newcommand{\yes}{\textbullet}
\newcommand{\no}{}

% ─── Document ────────────────────────────────────────────────────────────────
\begin{document}
\pagenumbering{Alph}  % avoid duplicate page.1 between titlepage and body

% ─── Title page ──────────────────────────────────────────────────────────────
\begin{titlepage}
\centering
\vspace*{3cm}

{\Huge\bfseries\sffamily C Shim Architecture}\\[0.5cm]
{\LARGE\sffamily Existence Proof and Functional Audit}\\[2cm]

\begin{tabular}{@{}r@{\quad}l@{}}
  \textbf{Date}       & 2026-03-09 \\[4pt]
  \textbf{Scope}      & \file{src/zenoh\_dart.\{h,c\}} --- 62 exported symbols, \raisebox{0.5ex}{\texttildelow}700 lines \\[4pt]
  \textbf{Reviewed against} & zenoh-c v1.7.2, Dart SDK 3.11.0 (FVM), ffigen v20.1.1 \\[4pt]
  \textbf{Audience}   & Expert C and Dart code reviewers \\[4pt]
  \textbf{Supersedes} & \file{CA-C\_shim\_audit.md}, \file{CA-C\_shim\_proof.md}, \file{CC-C\_shim\_audit.md} \\
\end{tabular}

\vfill

{\large\sffamily zenoh-dart Project}\\[0.5cm]
{\small All Dart API documentation references verified against Dart 3.11.1 (latest stable).}

\end{titlepage}

% ─── Table of contents ───────────────────────────────────────────────────────
\tableofcontents
\clearpage
\pagenumbering{arabic}  % restart at page 1 for body

% ══════════════════════════════════════════════════════════════════════════════
\part{Why the C Shim Must Exist}
\label{part:why}
% ══════════════════════════════════════════════════════════════════════════════

The zenoh-c v1.7.2 public API relies on six patterns that Dart FFI cannot
call directly. This part proves each pattern exists in the zenoh-c headers
with exact file references, explains why Dart FFI cannot handle it with
references to official Dart documentation, and shows how the C shim
resolves each one.

% ──────────────────────────────────────────────────────────────────────────────
\section{How Dart FFI Resolves Symbols}
\label{sec:ffi-symbols}
% ──────────────────────────────────────────────────────────────────────────────

\subsection{Dart FFI Is C-Only}

The \code{dart:ffi} library documentation (Dart 3.11.1) states:

\begin{quotebox}
``Foreign Function Interface for interoperability with the C programming
language.''\\
--- \href{https://api.dart.dev/stable/latest/dart-ffi/dart-ffi-library.html}{dart:ffi library overview}
\end{quotebox}

\begin{quotebox}
``Dart mobile, command-line, and server apps running on the Dart Native
platform can use the dart:ffi library to call native C APIs, and to read,
write, allocate, and deallocate native memory.''\\
--- \href{https://dart.dev/interop/c-interop}{C interop guide}
\end{quotebox}

The binding generator (\code{ffigen} v20.1.1) reinforces this constraint:

\begin{quotebox}
``Note: FFIgen only supports parsing \texttt{C} headers, not \texttt{C++} headers.''\\
--- \href{https://pub.dev/packages/ffigen}{ffigen package}
\end{quotebox}

\textbf{No Dart 3.x release has added support for calling C macros,
\code{static inline} functions, or C++ from \code{dart:ffi}.} The
\href{https://dart.dev/resources/whats-new}{Dart what's new} page for
releases 3.0 through 3.11 contains no FFI changes related to these
constructs. Dart ``build hooks'' (formerly ``native assets'') change how
native libraries are \emph{built and bundled} but do not alter how
\code{dart:ffi} resolves symbols at runtime.


\subsection{Symbol Resolution via \code{dlsym}}

\code{DynamicLibrary.open} loads a shared library and provides access to its
symbols:

\begin{quotebox}
``A dynamically loaded library is a mapping from symbols to memory addresses.''\\
--- \href{https://api.dart.dev/stable/latest/dart-ffi/DynamicLibrary-class.html}{DynamicLibrary class}
\end{quotebox}

\code{DynamicLibrary.lookup} resolves a symbol name to a memory address ---
functionally equivalent to POSIX \code{dlsym(3)}:

\begin{quotebox}
``Looks up a symbol in the DynamicLibrary and returns its address in memory.''\\
``Similar to the functionality of the dlsym(3) system call.''\\
``The symbol must be provided by the dynamic library.''\\
--- \href{https://api.dart.dev/stable/latest/dart-ffi/DynamicLibrary/lookup.html}{DynamicLibrary.lookup}
\end{quotebox}

\textbf{Consequence:} Only symbols present in the \code{.so}'s dynamic symbol
table (\code{nm -D}) are accessible. The following C constructs produce
\textbf{no symbols}:

\begin{table}[ht]
\centering
\begin{tabularx}{\textwidth}{lL}
  \toprule
  \textbf{Construct} & \textbf{Why No Symbol} \\
  \midrule
  \code{\#define} macros        & Text substitution by the preprocessor. No compiled code. \\
  \code{static inline} functions & Internal linkage + inlined. Not exported. \\
  \code{\_Generic} macros       & Compile-time type dispatch. Preprocessor construct. \\
  C++ overloaded functions       & Name-mangled. No stable C ABI symbol. \\
  \bottomrule
\end{tabularx}
\caption{C constructs that produce no dynamic symbols}
\label{tab:no-symbols}
\end{table}


\subsection{Confirmed by ffigen Bug Tracker}
\label{sec:ffigen-bugs}

The \code{ffigen} maintainers have explicitly addressed \code{static inline}
functions in two GitHub issues:

\textbf{Issue \#146} (``Unavailable inline C function''): When \code{ffigen}
previously generated bindings for inline functions like
\code{static inline uint16\_t libusb\_cpu\_to\_le16(...)}, the generated Dart
code referenced a dynamic library symbol that did not exist. The team's
resolution (merged in ffigen 2.0.1) was to \textbf{filter out inline functions}
entirely:

\begin{quotebox}
``We should skip them. (Generating an implementation for an arbitrary
function would be very involved.)''\\
--- \href{https://github.com/dart-lang/ffigen/issues/146}{dart-lang/ffigen\#146}, closed as COMPLETED
\end{quotebox}

\textbf{Issue \#459} (``Inline function available in compiled dylib, but no
bindings generated''): Even when functions are marked \code{extern inline}
(making them available in the \code{.so}), \code{ffigen} initially refused to
generate bindings. A partial fix (PR \#594) was merged to allow
\code{extern inline} functions, but \code{static inline} functions remain
\textbf{permanently excluded} --- they are not symbols, so there is nothing
to bind to.

\begin{quotebox}
``Static inline functions don't appear in compiled dylibs and cannot be
bound --- this behavior is expected.''\\
--- \href{https://github.com/dart-lang/native/issues/459}{dart-lang/native\#459}
\end{quotebox}

This confirms that \code{static inline} is a known, permanent limitation of
Dart FFI, not an oversight that might be fixed in a future Dart release.


\subsection{How the Shim Bridges This}

The C shim (\file{src/zenoh\_dart.c}) is compiled as a C translation unit
that \code{\#include}s the zenoh-c headers. During compilation, the C
compiler:

\begin{enumerate}
  \item Expands all \code{\#define} macros
  \item Inlines all \code{static inline} functions
  \item Resolves all \code{\_Generic} dispatches
\end{enumerate}

The shim's functions are marked \code{FFI\_PLUGIN\_EXPORT}
(resolving to \code{\_\_attribute\_\_((visi\-bi\-li\-ty("default")))}),
which ensures they appear in \code{libzenoh\_dart.so}'s dynamic symbol
table. Dart FFI can then call them via \code{DynamicLibrary.lookup}.


% ──────────────────────────────────────────────────────────────────────────────
\section{Pattern 1: \code{static inline} Move Functions}
\label{sec:pattern1}
% ──────────────────────────────────────────────────────────────────────────────

\subsection{The Problem}

zenoh-c uses an ownership model where consuming a resource requires
converting a \type{z\_owned\_\allowbreak{}*\_t*} pointer to a
\type{z\_moved\_\allowbreak{}*\_t*} pointer. This conversion is done via
\code{static inline} functions defined in
\file{extern/\allowbreak{}zenoh-c/\allowbreak{}include/\allowbreak{}zenoh\_macros.h}.

\textbf{There are 56 such functions.} Here are the ones used by the C shim:

\begin{lstlisting}[style=cstyle,caption={Move functions from \file{zenoh\_macros.h} (lines 7--51)}]
static inline z_moved_bytes_t* z_bytes_move(z_owned_bytes_t* x) {
    return (z_moved_bytes_t*)(x);
}
static inline z_moved_config_t* z_config_move(z_owned_config_t* x) {
    return (z_moved_config_t*)(x);
}
static inline z_moved_session_t* z_session_move(z_owned_session_t* x) {
    return (z_moved_session_t*)(x);
}
static inline z_moved_publisher_t* z_publisher_move(z_owned_publisher_t* x) {
    return (z_moved_publisher_t*)(x);
}
static inline z_moved_subscriber_t* z_subscriber_move(z_owned_subscriber_t* x) {
    return (z_moved_subscriber_t*)(x);
}
static inline z_moved_encoding_t* z_encoding_move(z_owned_encoding_t* x) {
    return (z_moved_encoding_t*)(x);
}
// ... 50 more variants for other types
\end{lstlisting}

These are pointer casts at the C level --- they reinterpret the same memory
address as a different pointer type. But they are \code{static inline}:

\begin{itemize}
  \item \code{static} --- internal linkage (not visible outside the
        translation unit)
  \item \code{inline} --- the compiler may substitute the function body at
        the call site
\end{itemize}

\textbf{Neither property produces an exported symbol.}


\subsection{Symbol Table Proof}

Running \code{nm -D} on the compiled \code{libzenohc.so} confirms these
symbols do not exist:

\begin{lstlisting}[style=shellstyle]
$ nm -D libzenohc.so | grep z_bytes_move
(no output)

$ nm -D libzenohc.so | grep z_config_move
(no output)

$ nm -D libzenohc.so | grep z_session_move
(no output)
\end{lstlisting}

Zero results. These functions exist only in the header file. Dart FFI's
\fn{DynamicLibrary.lookup('z\_bytes\_move')} would throw an
\code{ArgumentError}.


\subsection{Why These Are Required}

The move functions are not optional convenience wrappers. zenoh-c's core API
functions \textbf{require} \type{z\_moved\_*\_t*} parameters. Without the
move functions, these core functions cannot be called:

\begin{lstlisting}[style=cstyle,caption={Exported functions requiring moved parameters}]
// z_open REQUIRES z_moved_config_t* (line 3796)
ZENOHC_API z_result_t z_open(
    struct z_owned_session_t *this_,
    struct z_moved_config_t *config,        // <- requires move
    const struct z_open_options_t *_options);

// z_put REQUIRES z_moved_bytes_t* (line 4062)
ZENOHC_API z_result_t z_put(
    const struct z_loaned_session_t *session,
    const struct z_loaned_keyexpr_t *key_expr,
    struct z_moved_bytes_t *payload,         // <- requires move
    struct z_put_options_t *options);

// z_config_drop REQUIRES z_moved_config_t* (line 2214)
ZENOHC_API void z_config_drop(struct z_moved_config_t *this_);

// z_declare_subscriber REQUIRES z_moved_closure_sample_t* (line 2331)
ZENOHC_API z_result_t z_declare_subscriber(
    const struct z_loaned_session_t *session,
    struct z_owned_subscriber_t *subscriber,
    const struct z_loaned_keyexpr_t *key_expr,
    struct z_moved_closure_sample_t *callback,  // <- requires move
    struct z_subscriber_options_t *options);
\end{lstlisting}

These functions \textbf{are} exported (confirmed via \code{nm -D}):

\begin{lstlisting}[style=shellstyle]
$ nm -D libzenohc.so | grep " T z_open$"
000000000039acb0 T z_open

$ nm -D libzenohc.so | grep " T z_put$"
000000000039e3f0 T z_put

$ nm -D libzenohc.so | grep " T z_config_drop$"
000000000038f790 T z_config_drop
\end{lstlisting}

\begin{patternbox}
\textbf{The functions exist. Their required parameter types can only be
constructed via functions that do not exist.} This is the fundamental barrier.
\end{patternbox}


\subsection{What the Moved Types Actually Are}

The \type{z\_moved\_*\_t} types are wrapper structs defined in
\file{zenoh\_commons.h}:

\begin{lstlisting}[style=cstyle]
typedef struct z_moved_bytes_t {
    struct z_owned_bytes_t _this;
} z_moved_bytes_t;

typedef struct z_moved_config_t {
    struct z_owned_config_t _this;
} z_moved_config_t;
\end{lstlisting}

At the machine level, \type{z\_moved\_bytes\_t*} and
\type{z\_owned\_bytes\_t*} point to the same memory --- the wrapper struct
has the same layout as its single member. The \code{static inline} move
function is an identity cast resolved at compile time with zero runtime cost.

\textbf{Could Dart perform this cast directly?} In theory yes, but:

\begin{enumerate}
  \item Dart FFI has no mechanism to express the type distinction between
        \type{z\_owned\_bytes\_t*} and \type{z\_moved\_bytes\_t*} --- both
        would be \type{Pointer<Void>}.
  \item \code{ffigen} v20.1.1 \textbf{permanently excludes}
        \code{static inline} functions from binding generation
        (see \S\ref{sec:ffigen-bugs}).
  \item The type safety that the move functions provide (preventing
        accidental reuse of consumed resources) would be lost entirely.
\end{enumerate}


\subsection{How the C Shim Solves This}

The shim calls the \code{static inline} move functions at compile time and
exports the result as real functions:

\begin{lstlisting}[style=cstyle,caption={Shim resolves \fn{z\_config\_move} at compile time}]
FFI_PLUGIN_EXPORT void zd_config_drop(z_owned_config_t* config) {
    z_config_drop(z_config_move(config));
    //            ^^^^^^^^^^^^^^^^^^^^^^
    //            static inline, resolved at compile time
}

FFI_PLUGIN_EXPORT int zd_open_session(z_owned_session_t* session,
                                      z_owned_config_t* config) {
    return z_open(session, z_config_move(config), NULL);
}

FFI_PLUGIN_EXPORT int zd_put(
    const z_loaned_session_t* session,
    const z_loaned_keyexpr_t* keyexpr,
    z_owned_bytes_t* payload) {
    z_put_options_t opts;
    z_put_options_default(&opts);
    return z_put(session, keyexpr, z_bytes_move(payload), &opts);
}
\end{lstlisting}

The shim applies this pattern consistently across all 8 drop functions:

\begin{itemize}[nosep]
  \item \fn{zd\_config\_drop}, \fn{zd\_bytes\_drop}, \fn{zd\_string\_drop}
  \item \fn{zd\_subscriber\_drop}, \fn{zd\_publisher\_drop}, \fn{zd\_session\_drop}
  \item \fn{zd\_shm\_mut\_drop}, \fn{zd\_shm\_provider\_drop}
\end{itemize}


% ──────────────────────────────────────────────────────────────────────────────
\section{Pattern 2: C11 \code{\_Generic} Polymorphic Macros}
\label{sec:pattern2}
% ──────────────────────────────────────────────────────────────────────────────

\subsection{The Problem}

zenoh-c provides convenience macros that dispatch to the correct type-specific
function based on the argument type. These use C11's \code{\_Generic} keyword,
a compile-time type-selection mechanism (ISO~C11, \S6.5.1.1).

From \file{extern/zenoh-c/include/zenoh\_macros.h}:

\begin{lstlisting}[style=cstyle,caption={\code{z\_loan} macro (lines 65--119, 51 branches)}]
#define z_loan(this_) \
    _Generic((this_), \
        z_owned_bytes_t : z_bytes_loan, \
        z_owned_config_t : z_config_loan, \
        z_owned_session_t : z_session_loan, \
        z_owned_publisher_t : z_publisher_loan, \
        z_owned_subscriber_t : z_subscriber_loan, \
        /* ... 40+ more branches ... */ \
    )(&this_)
\end{lstlisting}

Four such macros exist: \code{z\_loan} (51~branches), \code{z\_drop}
(48~branches), \code{z\_move} (56~branches), and \code{z\_loan\_mut}
(25~branches). None produce symbols:

\begin{lstlisting}[style=shellstyle]
$ nm -D libzenohc.so | grep "z_loan$"
(no output)

$ nm -D libzenohc.so | grep "z_drop$"
(no output)

$ nm -D libzenohc.so | grep "z_move$"
(no output)
\end{lstlisting}


\subsection{Compounding With Pattern 1}

The underlying monomorphic functions (\fn{z\_bytes\_loan},
\fn{z\_config\_loan}, etc.) \textbf{are} exported. But the \code{z\_move}
macro dispatches to the \code{static inline} move functions (\pattern{1}).
So \code{z\_move(myConfig)} expands to \code{z\_config\_move(\&myConfig)},
which is \code{static inline} and not exported.

The C shim bypasses both layers:

\begin{lstlisting}[style=cstyle]
// Instead of: z_drop(z_move(config))
// The shim does:
z_config_drop(z_config_move(config));
// z_config_drop  -> exported real function (nm -D confirms)
// z_config_move  -> static inline, resolved at shim compile time
\end{lstlisting}


% ──────────────────────────────────────────────────────────────────────────────
\section{Pattern 3: Options Struct Initialization}
\label{sec:pattern3}
% ──────────────────────────────────────────────────────────────────────────────

\subsection{The Problem}

zenoh-c functions take options structs that must be initialized before use.
While the initialization functions (\fn{z\_put\_options\_default}, etc.) are
real exported functions, the barrier is the \textbf{combination} of:

\begin{enumerate}
  \item Allocating the options struct with the correct size
  \item Initializing it via the default function
  \item Setting fields that require move semantics
        (\code{opts.encoding = z\_encoding\_move(...)})
  \item Passing it to the target function
        (\code{z\_put(session, keyexpr, z\_bytes\_move(payload), \&opts)})
\end{enumerate}

Steps~3 and~4 require \code{static inline} functions (\pattern{1}).


\subsection{The Struct Layout Problem}

The \type{z\_put\_options\_t} struct has 8 fields with platform-dependent
alignment:

\begin{lstlisting}[style=cstyle,caption={\type{z\_put\_options\_t} layout from \file{zenoh\_commons.h}}]
typedef struct z_put_options_t {
    struct z_moved_encoding_t *encoding;
    z_congestion_control_t congestion_control;
    z_priority_t priority;
    bool is_express;
    struct z_moved_bytes_t *attachment;
    z_reliability_t reliability;
    const struct z_loaned_timestamp_t *timestamp;
    struct z_moved_source_info_t *source_info;
} z_put_options_t;
\end{lstlisting}

Dart FFI's
\href{https://api.dart.dev/stable/latest/dart-ffi/Struct-class.html}{Struct class}
requires explicit field annotations (\code{@Int32()}, \code{@Pointer()})
to access struct fields. The \file{ffigen.yaml} maps all zenoh-c types to
\href{https://api.dart.dev/stable/latest/dart-ffi/Opaque-class.html}{Opaque},
not \code{Struct}, because the layout changes between versions.


\subsection{How the C Shim Solves This}

The shim internalizes the entire options workflow. For functions needing
selective overrides, it uses sentinel parameters (\code{NULL} for strings,
\code{-1} for enums):

\begin{lstlisting}[style=cstyle,caption={Sentinel-based optional parameters in \fn{zd\_declare\_publisher}}]
FFI_PLUGIN_EXPORT int zd_declare_publisher(
    const z_loaned_session_t* session,
    z_owned_publisher_t* publisher,
    const z_loaned_keyexpr_t* keyexpr,
    const char* encoding,         // NULL = default
    int congestion_control,       // -1 = default
    int priority) {               // -1 = default
    z_publisher_options_t opts;
    z_publisher_options_default(&opts);

    z_owned_encoding_t owned_encoding;
    if (encoding != NULL) {
        z_encoding_from_str(&owned_encoding, encoding);
        opts.encoding = z_encoding_move(&owned_encoding);
    }
    if (congestion_control >= 0) {
        opts.congestion_control = (z_congestion_control_t)congestion_control;
    }
    if (priority >= 0) {
        opts.priority = (z_priority_t)priority;
    }

    return z_declare_publisher(session, publisher, keyexpr, &opts);
}
\end{lstlisting}

\textbf{Dart sees simple scalar parameters:}

\begin{lstlisting}[style=dartstyle,caption={Dart invocation --- no options struct, no move functions}]
final rc = bindings.zd_declare_publisher(
    loanedSession.cast(),
    ptr.cast(),
    loanedKe.cast(),
    encodingStr.cast(),           // Pointer<Utf8> or nullptr
    congestionControl.index,      // int (0 or 1)
    priority.index + 1,           // int (1-7)
);
\end{lstlisting}


% ──────────────────────────────────────────────────────────────────────────────
\section{Pattern 4: Opaque Type Sizes}
\label{sec:pattern4}
% ──────────────────────────────────────────────────────────────────────────────

\subsection{The Problem}

zenoh-c types are opaque structs whose sizes are determined at zenoh-c
compile time, injected via cbindgen:

\begin{lstlisting}[style=cstyle,caption={Opaque type definitions from \file{zenoh\_opaque.h}}]
typedef struct ALIGN(8) z_owned_config_t {
  uint8_t _0[2008];       // 2008 bytes
} z_owned_config_t;

typedef struct ALIGN(8) z_owned_session_t {
  uint8_t _0[8];          // 8 bytes
} z_owned_session_t;

typedef struct ALIGN(8) z_owned_bytes_t {
  uint8_t _0[40];         // 40 bytes
} z_owned_bytes_t;
\end{lstlisting}

Dart FFI's \href{https://api.dart.dev/stable/latest/dart-ffi/sizeOf.html}{\code{sizeOf<T>()}}
only works for types whose layout is known at Dart compile time. For
\code{Opaque} types, \code{sizeOf} is not defined.


\subsection{How the C Shim Solves This}

The shim exports \code{sizeof} queries as real functions (10 total):

\begin{lstlisting}[style=cstyle]
FFI_PLUGIN_EXPORT size_t zd_config_sizeof(void) {
    return sizeof(z_owned_config_t);   // resolves to 2008
}
FFI_PLUGIN_EXPORT size_t zd_session_sizeof(void) {
    return sizeof(z_owned_session_t);  // resolves to 8
}
FFI_PLUGIN_EXPORT size_t zd_bytes_sizeof(void) {
    return sizeof(z_owned_bytes_t);    // resolves to 40
}
\end{lstlisting}

Every Dart wrapper class queries the size at construction time:

\begin{lstlisting}[style=dartstyle]
final size = bindings.zd_session_sizeof();    // query at runtime
final Pointer<Void> ptr = calloc.allocate(size);  // allocate exactly
\end{lstlisting}

This pattern is version-safe: if zenoh-c changes struct sizes in a future
release, the Dart code automatically allocates the correct amount.


% ──────────────────────────────────────────────────────────────────────────────
\section{Pattern 5: Closure Callbacks Across Thread Boundaries}
\label{sec:pattern5}
% ──────────────────────────────────────────────────────────────────────────────

\subsection{The Problem}

zenoh-c delivers asynchronous events via closure structs containing function
pointers and a \code{void*} context:

\begin{lstlisting}[style=cstyle]
typedef struct z_owned_closure_sample_t {
  void *_context;
  void (*_call)(struct z_loaned_sample_t *sample, void *context);
  void (*_drop)(void *context);
} z_owned_closure_sample_t;
\end{lstlisting}

The \fn{\_call} function pointer is invoked by zenoh-c's internal Tokio
runtime threads, not by the Dart isolate's thread.


\subsection{Why Dart FFI Cannot Construct These Closures}

There are three independent blocking constraints:

\subsubsection{Constraint A: No closures in \code{Pointer.fromFunction}}

\begin{quotebox}
``Creates a Dart function pointer from a top-level or static Dart function.
{[...]} The function must not be a closure (i.e.\ it must not capture any
local variables).''\\
--- \href{https://api.dart.dev/stable/latest/dart-ffi/Pointer/fromFunction.html}{Pointer.fromFunction docs}
\end{quotebox}

Dart cannot create a C function pointer that carries state (like which
\code{StreamController} to forward samples to).

\subsubsection{Constraint B: Non-Dart thread invocation is undefined behavior}

\href{https://github.com/dart-lang/sdk/issues/48865}{dart-lang/sdk\#48865}
confirms that calling a \fn{Pointer.fromFunction} callback from a non-Dart
thread causes undefined behavior.

\subsubsection{Constraint C: \code{NativeCallable.listener} limitations}

Dart 3.1+ added \code{NativeCallable.listener} which can be called from any
thread, but it still cannot capture closures, and context must be passed via
integer baton rather than \code{void*} pointer.
See \href{https://github.com/dart-lang/sdk/issues/52689}{dart-lang/sdk\#52689}.


\subsection{How the C Shim Solves This}

The shim implements a \textbf{NativePort callback bridge}:

\begin{enumerate}
  \item \textbf{C side:} A static C callback extracts fields from the
        zenoh-c sample, constructs a \code{Dart\_CObject} array, and posts
        it to a Dart \code{ReceivePort} via the thread-safe function
        \fn{Dart\_\allowbreak{}Post\-CObject\_DL}.
  \item \textbf{Dart side:} A \code{ReceivePort} listener on the Dart event
        loop receives the deserialized data and forwards it to a
        \code{StreamController<Sample>}.
\end{enumerate}

\begin{lstlisting}[style=cstyle,caption={NativePort callback bridge (C side)}]
typedef struct {
  Dart_Port_DL dart_port;
} zd_subscriber_context_t;

static void _zd_sample_callback(z_loaned_sample_t* sample, void* context) {
  zd_subscriber_context_t* ctx = (zd_subscriber_context_t*)context;
  // Extract fields, build Dart_CObject array...
  Dart_PostCObject_DL(ctx->dart_port, &c_array);  // thread-safe
}

FFI_PLUGIN_EXPORT int zd_declare_subscriber(
    const z_loaned_session_t* session,
    z_owned_subscriber_t* subscriber,
    const z_loaned_keyexpr_t* keyexpr,
    int64_t dart_port) {  // receives port number, not function pointer
  zd_subscriber_context_t* ctx = malloc(sizeof(zd_subscriber_context_t));
  ctx->dart_port = (Dart_Port_DL)dart_port;
  z_owned_closure_sample_t callback;
  z_closure_sample(&callback, _zd_sample_callback, _zd_sample_drop, ctx);
  return z_declare_subscriber(session, subscriber, keyexpr,
                              z_closure_sample_move(&callback), NULL);
}
\end{lstlisting}

\fn{Dart\_Post\-CObject\_DL} is
\href{https://api.dart.dev/stable/latest/dart-ffi/NativeApi/initializeApiDLData.html}{documented as safe}
to call from any thread, performs deep copy semantics, and requires
one-time initialization via \fn{Dart\_\allowbreak{}InitializeApiDL}.

The same pattern is used for all three callback types:

\begin{table}[ht]
\centering\footnotesize
\begin{tabularx}{\textwidth}{lLlL}
  \toprule
  \textbf{Callback} & \textbf{C function} & \textbf{Data posted} & \textbf{Dart receiver} \\
  \midrule
  Subscriber   & \fn{\_zd\_sample\_callback}           & 5-element array & \code{Subscriber.stream} \\
  Matching     & \fn{\_zd\_matching\_status\_callback}  & Int64 (0/1)  & \code{Publisher.matchingStatus} \\
  Scout        & \fn{\_zd\_scout\_hello\_callback}      & 3-element array & \code{Zenoh.scout()} future \\
  \bottomrule
\end{tabularx}
\caption{NativePort callback bridge usage across the shim}
\label{tab:callbacks}
\end{table}


% ──────────────────────────────────────────────────────────────────────────────
\section{Pattern 6: Loaning and Const/Mut Enforcement}
\label{sec:pattern6}
% ──────────────────────────────────────────────────────────────────────────────

\subsection{The Problem}

zenoh-c implements a \textbf{borrowing protocol} where owned types
(\code{z\_owned\_*\_t}) must be ``loaned'' to produce temporary
references (\code{z\_loaned\_*\_t}) before they can be passed to
most API functions. The protocol distinguishes immutable borrows
(\code{const z\_loaned\_*\_t*}) from mutable borrows
(\code{z\_loaned\_*\_t*}).

\begin{lstlisting}[style=cstyle,caption={Loan function pair from zenoh\_commons.h}]
// zenoh_commons.h:2218-2222
ZENOHC_API const struct z_loaned_config_t *z_config_loan(
    const struct z_owned_config_t *this_);              // immutable borrow
ZENOHC_API struct z_loaned_config_t *z_config_loan_mut(
    struct z_owned_config_t *this_);                    // mutable borrow
\end{lstlisting}

\subsection{Why the Per-Type Functions ARE Exported}

Unlike Pattern~1 (\code{static inline} move functions), the per-type loan
functions \textbf{are real, exported symbols} in \file{libzenohc.so}:

\begin{lstlisting}[style=shellstyle,basicstyle=\ttfamily\footnotesize\color{shellForeground}]
$ nm -D libzenohc.so | grep -E 'z_(config|session|bytes|publisher)_loan'
000000000038ad40 T z_bytes_loan
000000000038ad50 T z_bytes_loan_mut
0000000000390370 T z_config_loan
0000000000390380 T z_config_loan_mut
000000000039e810 T z_publisher_loan
00000000003a59b0 T z_session_loan
\end{lstlisting}

This is critical: these functions could theoretically be called from Dart
FFI via \fn{DynamicLibrary.lookup}. However, two independent barriers
prevent direct use.

\subsection{Barrier A: Macro Dispatch}

The zenoh-c API presents loaning through \code{\_Generic} macros
(\pattern{2}). No symbol named \code{z\_loan} or \code{z\_loan\_mut}
exists in \file{libzenohc.so}. Dart must know which monomorphic function
to call for each type---a responsibility the shim absorbs.

\subsection{Barrier B: Const Qualification Erasure}

Even if Dart called the per-type functions directly, Dart FFI's type
system \textbf{erases the const/mut distinction}. All zenoh-c types are
mapped to \code{ffi.Opaque} in \file{ffigen.yaml}. A
\code{Pointer<z\_loaned\_config\_t>} in Dart carries no information
about whether the pointer is \code{const} or mutable.

\begin{quotebox}
``A Pointer represents a pointer into native C memory.''\\
--- \href{https://api.dart.dev/stable/latest/dart-ffi/Pointer-class.html}{Pointer class} (Dart 3.11.1)
\end{quotebox}

There is no \code{ConstPointer<T>} in \code{dart:ffi}. This means:

\begin{itemize}
  \item Dart cannot distinguish \code{const z\_loaned\_config\_t*} from
        \code{z\_loaned\_config\_t*}
  \item There is no compile-time or runtime check preventing mutable
        operations on an immutably-borrowed reference
  \item Incorrect usage would cause undefined behavior in zenoh-c
\end{itemize}

\subsection{How the Shim Resolves This}

\begin{enumerate}
  \item \textbf{Macro bypass:} Calls monomorphic functions directly
        (\fn{z\_config\_loan} not \code{z\_loan(\allowbreak{}config)})
  \item \textbf{Const enforcement in C:} Function signatures enforce
        const correctness at the C compiler level
  \item \textbf{Naming convention:} \code{zd\_*\_loan} for immutable
        borrows, \code{zd\_*\_loan\_mut} for mutable borrows
\end{enumerate}

\begin{lstlisting}[style=cstyle,caption={Loan wrappers enforce const at the C level}]
// Returns const pointer -- C compiler prevents mutation
FFI_PLUGIN_EXPORT const z_loaned_config_t*
    zd_config_loan(const z_owned_config_t* config);

// Returns mutable pointer -- allows mutation (SHM buffer writes)
FFI_PLUGIN_EXPORT z_loaned_shm_mut_t*
    zd_shm_mut_loan_mut(z_owned_shm_mut_t* buf);
\end{lstlisting}

\subsection{Why This Is Distinct From Patterns 1 and 2}

\begin{table}[ht]
\centering\footnotesize
\begin{tabularx}{\textwidth}{lLLL}
  \toprule
  \textbf{Aspect} & \textbf{Pattern 1} & \textbf{Pattern 2} & \textbf{Pattern 6} \\
  \midrule
  Barrier & No exported symbol (\code{static inline}) & No exported symbol (\code{\_Generic} macro) & Exported symbol exists, but const/mut qualifier is lost across FFI \\[4pt]
  \code{nm -D} & Function absent & Macro absent & \textbf{Function present} (\code{T} in symbol table) \\[4pt]
  Could Dart call it? & No & No & Technically yes, but \textbf{unsafely} \\[4pt]
  What the shim adds & Inlines the call into exported body & Calls monomorphic function directly & Enforces const/mut at C level + provides stable naming \\
  \bottomrule
\end{tabularx}
\caption{Pattern~6 is a semantic barrier, not a syntactic one}
\label{tab:p6-vs-p1-p2}
\end{table}


% ──────────────────────────────────────────────────────────────────────────────
\section{End-to-End Worked Example: \code{z\_put}}
\label{sec:worked-example}
% ──────────────────────────────────────────────────────────────────────────────

This section traces a single operation --- publishing data on a key expression
--- through the six patterns to show how they compound.

\subsection{What a C Consumer Does}

\begin{lstlisting}[style=cstyle]
z_put_options_t opts;
z_put_options_default(&opts);                     // Pattern 3

z_owned_bytes_t payload;
z_bytes_copy_from_str(&payload, "Hello");

z_put(z_session_loan(&session),                   // Pattern 6 + Pattern 2
      z_view_keyexpr_loan(&keyexpr),              // Pattern 6 + Pattern 2
      z_bytes_move(&payload),                     // Pattern 1 (static inline)
      &opts);                                     // Pattern 3
\end{lstlisting}

\subsection{What Dart FFI Cannot Do}

Dart FFI \textbf{can} call \fn{z\_put} (exported),
\fn{z\_put\_\allowbreak{}options\_default} (exported),
and \fn{z\_bytes\_\allowbreak{}copy\_from\_str} (exported).

Dart FFI \textbf{cannot} call:
\fn{z\_bytes\_move} (\code{static inline}, \pattern{1}),
\code{sizeof(z\_put\_\allowbreak{}options\_t)} (\pattern{4}),
or set \code{opts.\allowbreak{}encoding}
(requires \fn{z\_encoding\_move}, \pattern{3}).

\subsection{The Full Call Chain}

\begin{lstlisting}[style=shellstyle,basicstyle=\ttfamily\footnotesize\color{shellForeground}]
Dart: session.put('demo/test', 'Hello')
  -> ZBytes.fromString('Hello')
    -> calloc.allocate(bindings.zd_bytes_sizeof())      [Pattern 4]
    -> bindings.zd_bytes_copy_from_str(ptr, nativeStr)   [exported]
  -> bindings.zd_session_loan(sessionPtr)                  [Pattern 6]
  -> bindings.zd_view_keyexpr_loan(kePtr)                  [Pattern 6]
  -> bindings.zd_put(loanedSession, loanedKe, payloadPtr)  [shim]
    -> C: z_put_options_default(&opts)                   [Pattern 3]
    -> C: z_put(session, keyexpr,
                z_bytes_move(payload),                   [Pattern 1]
                &opts)                                   [Pattern 3]
      -> libzenohc.so: z_put()                           [exported]
  -> payload.markConsumed()                              [Dart ownership]
\end{lstlisting}

\begin{patternbox}
Four patterns, one operation. The C shim collapses all of them into a single
FFI-callable function with simple scalar parameters.
\end{patternbox}


\subsection{What the Shim Does NOT Do}

The shim is deliberately thin. It does not add error recovery, buffer data,
manage threads, add abstraction layers, or change zenoh-c's ownership
semantics. Every \code{zd\_*} function maps to one or a small sequence of
\code{z\_*} calls. The shim is a \textbf{mechanical flattening layer}, not
an abstraction.


% ══════════════════════════════════════════════════════════════════════════════
\clearpage
\part{Function-by-Function Analysis}
\label{part:functions}
% ══════════════════════════════════════════════════════════════════════════════

% ──────────────────────────────────────────────────────────────────────────────
\section{Initialization (2 functions)}
% ──────────────────────────────────────────────────────────────────────────────

\begin{table}[ht]
\centering\small
\begin{tabularx}{\textwidth}{llL}
  \toprule
  \textbf{Function} & \textbf{Wraps} & \textbf{Purpose} \\
  \midrule
  \fn{zd\_init\_dart\_api\_dl} & \fn{Dart\_InitializeApiDL} & Initializes Dart native API for \fn{Dart\_PostCObject\_DL} \\
  \fn{zd\_init\_log} & \fn{zc\_init\_log\_from\_env\_or} & Initializes zenoh logger \\
  \bottomrule
\end{tabularx}
\end{table}

\textbf{Assessment:} Correct and minimal. Called once during
\fn{\_initBindings()}.


% ──────────────────────────────────────────────────────────────────────────────
\section{Config (5 functions)}
% ──────────────────────────────────────────────────────────────────────────────

\begin{table}[ht]
\centering\small
\begin{tabularx}{\textwidth}{llL}
  \toprule
  \textbf{Function} & \textbf{Wraps} & \textbf{Patterns} \\
  \midrule
  \fn{zd\_config\_sizeof} & \code{sizeof(z\_owned\_config\_t)} & P4 \\
  \fn{zd\_config\_default} & \fn{z\_config\_default} & P3 \\
  \fn{zd\_config\_insert\_json5} & \fn{z\_config\_loan\_mut} + \fn{zc\_config\_insert\_json5} & P1, P2 \\
  \fn{zd\_config\_loan} & \fn{z\_config\_loan} & P1, P2 \\
  \fn{zd\_config\_drop} & \fn{z\_config\_drop(z\_config\_move(...))} & P1 \\
  \bottomrule
\end{tabularx}
\end{table}


% ──────────────────────────────────────────────────────────────────────────────
\section{Session (4 functions)}
% ──────────────────────────────────────────────────────────────────────────────

\begin{table}[ht]
\centering\small
\begin{tabularx}{\textwidth}{llL}
  \toprule
  \textbf{Function} & \textbf{Wraps} & \textbf{Patterns} \\
  \midrule
  \fn{zd\_session\_sizeof} & \code{sizeof(z\_owned\_session\_t)} & P4 \\
  \fn{zd\_open\_session} & \fn{z\_open(session, z\_config\_move(config), NULL)} & P1, P3 \\
  \fn{zd\_session\_loan} & \fn{z\_session\_loan} & P1, P2 \\
  \fn{zd\_close\_session} & \fn{z\_close} + \fn{z\_session\_drop} & P1 \\
  \bottomrule
\end{tabularx}
\end{table}

\begin{findingbox}[Finding F1: \code{z\_close} return value ignored]{findingLow}
The return value of \fn{z\_close} is discarded in \fn{zd\_close\_session}.
zenoh-cpp similarly ignores it. Close-time errors are non-actionable.
\textbf{Severity: Low.}
\end{findingbox}


% ──────────────────────────────────────────────────────────────────────────────
\section{KeyExpr (4), Bytes (6), String/View String (7)}
% ──────────────────────────────────────────────────────────────────────────────

All are straightforward \pattern{1} + \pattern{4} wrappers. Key points:

\begin{itemize}
  \item View key expressions borrow the backing string; lifetime is
        guaranteed by \fn{\_withKeyExpr}'s \code{try/finally}.
  \item \fn{zd\_string\_data} returns non-null-terminated data; Dart
        correctly uses \code{toDart\-String(\allowbreak{}length:\ len)} everywhere.
  \item View string convenience functions (\fn{zd\_view\_string\_data/len})
        combine loan + access in one FFI round-trip.
\end{itemize}


% ──────────────────────────────────────────────────────────────────────────────
\section{Put / Delete (2 functions)}
% ──────────────────────────────────────────────────────────────────────────────

Both initialize options to defaults and pass through. \fn{zd\_put} moves the
payload; Dart calls \code{payload.markConsumed()} after. Session-level
\fn{zd\_put} intentionally omits encoding/QoS/attachment --- those are
available via \fn{zd\_publisher\_put}.


% ──────────────────────────────────────────────────────────────────────────────
\section{Subscriber (3 functions + 2 static callbacks)}
% ──────────────────────────────────────────────────────────────────────────────

Resolves Patterns 1, 4, and 5. The subscriber callback is the primary
motivation for \pattern{5}. Context is heap-allocated via \code{malloc},
freed by \fn{\_zd\_sample\_drop} when the subscriber is undeclared.
Error-path closure drop is defensively correct.


% ──────────────────────────────────────────────────────────────────────────────
\section{Publisher (8 functions + 2 static callbacks)}
% ──────────────────────────────────────────────────────────────────────────────

Resolves all five patterns. Both payload and attachment are moved
(consumed) in \fn{zd\_pub\-lisher\_put}. The matching listener uses
the NativePort bridge (\pattern{5}).


% ──────────────────────────────────────────────────────────────────────────────
\section{Info / ZID (4 functions)}
% ──────────────────────────────────────────────────────────────────────────────

The two ZID collection functions
(\fn{zd\_info\_routers\_zid}, \fn{zd\_info\_peers\_zid})
use synchronous buffer-based collection (stack-allocated context,
\code{NULL} drop callback) instead of NativePort---appropriate because
these are synchronous. Buffer holds up to 64~ZIDs (1024~bytes).


% ──────────────────────────────────────────────────────────────────────────────
\section{Scout (2 functions + 1 static callback)}
% ──────────────────────────────────────────────────────────────────────────────

Uses NativePort with \textbf{stack-allocated} context (safe because
\fn{z\_scout} blocks until timeout). Posts null sentinel on completion;
Dart \code{Completer} resolves on null.

\begin{findingbox}[Finding F6: Scout return value ignored]{findingMedium}
\fn{Zenoh.scout()} does not check the return value from \fn{zd\_scout}.
Scout failures are indistinguishable from ``no entities found.''
\textbf{Severity: Medium.}
\end{findingbox}


% ──────────────────────────────────────────────────────────────────────────────
\section{Shared Memory (13 functions, conditionally compiled)}
% ──────────────────────────────────────────────────────────────────────────────

All guarded by:

\smallskip
\code{\#if defined(Z\_FEATURE\_SHARED\_MEMORY)}\\
\code{\ \ \ \ \&\& defined(Z\_FEATURE\_UNSTABLE\_API)}
\smallskip

\noindent
Alloc results are flattened from a compound result struct to return code +
out-parameter. Dart returns \code{null} on allocation failure (not throwing).


% ══════════════════════════════════════════════════════════════════════════════
\clearpage
\part{Cross-Cutting Analysis}
\label{part:cross-cutting}
% ══════════════════════════════════════════════════════════════════════════════

% ──────────────────────────────────────────────────────────────────────────────
\section{Memory Ownership Protocol}
% ──────────────────────────────────────────────────────────────────────────────

The Dart side implements a three-state model mirroring zenoh-c's gravestone
semantics:

\begin{table}[ht]
\centering
\begin{tabular}{lll}
  \toprule
  \textbf{State} & \textbf{Flags} & \textbf{Operations} \\
  \midrule
  Live      & \code{\_disposed=false, \_consumed=false} & All operations \\
  Consumed  & \code{\_consumed=true}                     & None (\code{StateError}) \\
  Disposed  & \code{\_disposed=true}                     & \code{dispose()} (no-op) \\
  \bottomrule
\end{tabular}
\end{table}

Two-step cleanup is applied consistently: (1)~\fn{zd\_*\_drop()} releases
zenoh-c internals (Rust allocator), then (2)~\code{calloc.free()} releases
the Dart-allocated wrapper memory.


% ──────────────────────────────────────────────────────────────────────────────
\section{Thread Safety}
% ──────────────────────────────────────────────────────────────────────────────

All C callbacks use only stack-local variables + read-only context access
+ \fn{Dart\_Post\-CObject\_DL} (thread-safe). No shared mutable state.
Context structs contain only a \type{Dart\_Port\_DL} integer, written
once during declaration and read from callbacks---no races possible.


% ──────────────────────────────────────────────────────────────────────────────
\section{Findings}
% ──────────────────────────────────────────────────────────────────────────────

\begin{findingbox}[F1: \code{z\_close} return value ignored]{findingLow}
\textbf{Severity: Low.} zenoh-cpp similarly ignores it. Close-time errors
are non-actionable.
\end{findingbox}

\begin{findingbox}[F2: Payload passes through \code{z\_bytes\_to\_string}]{findingLow}
\textbf{Severity: Low.} An unnecessary copy in the subscriber callback.
Sending raw bytes directly from \fn{z\_sample\_payload} would be more
efficient but requires the more complex bytes reader API.
\end{findingbox}

\begin{findingbox}[F3: \code{ZBytes.fromUint8List} element-by-element copy]{findingLow}
\textbf{Severity: Low.} Could use
\code{nativeBuf.asTypedList(len).setAll(0, data)} for bulk copy via
\code{memcpy}. The SHM path already uses this pattern correctly.
\end{findingbox}

\begin{findingbox}[F4: \code{z\_encoding\_from\_str} return code unchecked]{findingLow}
\textbf{Severity: Low.} In practice, zenoh-c's \fn{z\_encoding\_from\_str}
always succeeds (treats the string as an opaque MIME type). Defensive fix
would be to check the return code.
\end{findingbox}

\begin{findingbox}[F5: Allocator mismatch in \code{Zenoh.initLog}]{findingLow}
\textbf{Severity: Low.} \code{toNativeUtf8()} allocates via \code{malloc}
but cleanup uses \code{calloc.free()}. Both resolve to the same system
\code{free()} in practice, but violates the matching-allocator convention.
\end{findingbox}

\begin{findingbox}[F6: Scout return value ignored]{findingMedium}
\textbf{Severity: Medium.} \fn{Zenoh.scout()} does not check the return
value from \fn{zd\_scout}. Scout failures are indistinguishable from ``no
entities found.''
\end{findingbox}

\begin{findingbox}[F7: \code{ZBytes.dispose()} leaks calloc wrapper when consumed]{findingLow}
\textbf{Severity: Low.} When consumed, \code{dispose()} returns early
without calling \code{calloc.free(\_ptr)}. The native content was consumed
by zenoh-c, but the calloc wrapper still needs freeing.
\code{ShmMutBuffer.dispose()} handles this correctly.
\end{findingbox}


% ══════════════════════════════════════════════════════════════════════════════
\clearpage
\part{Pattern Resolution Matrix}
\label{part:matrix}
% ══════════════════════════════════════════════════════════════════════════════

This matrix maps every C shim function to which patterns it resolves.

{\small
\begin{longtable}{l Y{7mm} Y{7mm} Y{7mm} Y{7mm} Y{7mm} Y{7mm}}
  \toprule
  \textbf{C shim function} & \textbf{P1} & \textbf{P2} & \textbf{P3} & \textbf{P4} & \textbf{P5} & \textbf{P6} \\
  \midrule
  \endhead
  \bottomrule
  \endfoot
  \fn{zd\_init\_dart\_api\_dl}       & \no  & \no  & \no  & \no  & \yes & \no  \\
  \fn{zd\_init\_log}                 & \no  & \no  & \no  & \no  & \no  & \no  \\
  \fn{zd\_config\_sizeof}            & \no  & \no  & \no  & \yes & \no  & \no  \\
  \fn{zd\_config\_default}           & \no  & \no  & \yes & \no  & \no  & \no  \\
  \fn{zd\_config\_insert\_json5}     & \yes & \yes & \no  & \no  & \no  & \no  \\
  \fn{zd\_config\_loan}              & \no  & \yes & \no  & \no  & \no  & \yes \\
  \fn{zd\_config\_drop}              & \yes & \no  & \no  & \no  & \no  & \no  \\
  \fn{zd\_session\_sizeof}           & \no  & \no  & \no  & \yes & \no  & \no  \\
  \fn{zd\_open\_session}             & \yes & \no  & \yes & \no  & \no  & \no  \\
  \fn{zd\_session\_loan}             & \no  & \yes & \no  & \no  & \no  & \yes \\
  \fn{zd\_close\_session}            & \yes & \yes & \no  & \no  & \no  & \no  \\
  \fn{zd\_view\_keyexpr\_sizeof}     & \no  & \no  & \no  & \yes & \no  & \no  \\
  \fn{zd\_view\_keyexpr\_from\_str}  & \yes & \no  & \no  & \no  & \no  & \no  \\
  \fn{zd\_view\_keyexpr\_loan}       & \no  & \yes & \no  & \no  & \no  & \yes \\
  \fn{zd\_keyexpr\_as\_view\_string} & \yes & \no  & \no  & \no  & \no  & \no  \\
  \fn{zd\_bytes\_sizeof}             & \no  & \no  & \no  & \yes & \no  & \no  \\
  \fn{zd\_bytes\_copy\_from\_str}    & \yes & \no  & \no  & \no  & \no  & \no  \\
  \fn{zd\_bytes\_copy\_from\_buf}    & \yes & \no  & \no  & \no  & \no  & \no  \\
  \fn{zd\_bytes\_to\_string}         & \yes & \no  & \no  & \no  & \no  & \no  \\
  \fn{zd\_bytes\_loan}               & \no  & \yes & \no  & \no  & \no  & \yes \\
  \fn{zd\_bytes\_drop}               & \yes & \no  & \no  & \no  & \no  & \no  \\
  \fn{zd\_string\_sizeof}            & \no  & \no  & \no  & \yes & \no  & \no  \\
  \fn{zd\_string\_loan}              & \no  & \yes & \no  & \no  & \no  & \yes \\
  \fn{zd\_string\_data}              & \yes & \no  & \no  & \no  & \no  & \no  \\
  \fn{zd\_string\_len}               & \yes & \no  & \no  & \no  & \no  & \no  \\
  \fn{zd\_string\_drop}              & \yes & \no  & \no  & \no  & \no  & \no  \\
  \fn{zd\_view\_string\_sizeof}      & \no  & \no  & \no  & \yes & \no  & \no  \\
  \fn{zd\_view\_string\_data}        & \yes & \yes & \no  & \no  & \no  & \no  \\
  \fn{zd\_view\_string\_len}         & \yes & \yes & \no  & \no  & \no  & \no  \\
  \fn{zd\_put}                       & \yes & \no  & \yes & \no  & \no  & \no  \\
  \fn{zd\_delete}                    & \no  & \no  & \yes & \no  & \no  & \no  \\
  \fn{zd\_subscriber\_sizeof}        & \no  & \no  & \no  & \yes & \no  & \no  \\
  \fn{zd\_declare\_subscriber}       & \yes & \no  & \no  & \no  & \yes & \no  \\
  \fn{zd\_subscriber\_drop}          & \yes & \no  & \no  & \no  & \no  & \no  \\
  \fn{zd\_publisher\_sizeof}         & \no  & \no  & \no  & \yes & \no  & \no  \\
  \fn{zd\_declare\_publisher}        & \yes & \no  & \yes & \no  & \no  & \no  \\
  \fn{zd\_publisher\_loan}           & \no  & \yes & \no  & \no  & \no  & \yes \\
  \fn{zd\_publisher\_drop}           & \yes & \no  & \no  & \no  & \no  & \no  \\
  \fn{zd\_publisher\_put}            & \yes & \no  & \yes & \no  & \no  & \no  \\
  \fn{zd\_publisher\_delete}         & \no  & \no  & \yes & \no  & \no  & \no  \\
  \fn{zd\_publisher\_keyexpr}        & \yes & \no  & \no  & \no  & \no  & \no  \\
  \fn{zd\_pub..\_matching\_listener} & \yes & \no  & \no  & \no  & \yes & \no  \\
  \fn{zd\_pub..\_matching\_status}   & \no  & \no  & \no  & \no  & \no  & \no  \\
  \fn{zd\_info\_zid}                 & \yes & \no  & \no  & \no  & \no  & \no  \\
  \fn{zd\_id\_to\_string}            & \no  & \no  & \no  & \no  & \no  & \no  \\
  \fn{zd\_info\_routers\_zid}        & \no  & \no  & \no  & \no  & \yes & \no  \\
  \fn{zd\_info\_peers\_zid}          & \no  & \no  & \no  & \no  & \yes & \no  \\
  \fn{zd\_scout}                     & \yes & \no  & \yes & \no  & \yes & \no  \\
  \fn{zd\_whatami\_to\_view\_string} & \no  & \no  & \no  & \no  & \no  & \no  \\
  \fn{zd\_shm\_provider\_sizeof}     & \no  & \no  & \no  & \yes & \no  & \no  \\
  \fn{zd\_shm\_provider\_new}        & \yes & \no  & \no  & \no  & \no  & \no  \\
  \fn{zd\_shm\_provider\_loan}       & \no  & \yes & \no  & \no  & \no  & \yes \\
  \fn{zd\_shm\_provider\_drop}       & \yes & \no  & \no  & \no  & \no  & \no  \\
  \fn{zd\_shm\_provider\_available}  & \no  & \no  & \no  & \no  & \no  & \no  \\
  \fn{zd\_shm\_mut\_sizeof}          & \no  & \no  & \no  & \yes & \no  & \no  \\
  \fn{zd\_shm\_provider\_alloc}      & \yes & \no  & \no  & \no  & \no  & \no  \\
  \fn{zd\_shm..alloc\_gc\_defrag..}  & \yes & \no  & \no  & \no  & \no  & \no  \\
  \fn{zd\_shm\_mut\_loan\_mut}       & \no  & \yes & \no  & \no  & \no  & \yes \\
  \fn{zd\_shm\_mut\_data\_mut}       & \no  & \no  & \no  & \no  & \no  & \no  \\
  \fn{zd\_shm\_mut\_len}             & \no  & \no  & \no  & \no  & \no  & \no  \\
  \fn{zd\_bytes\_from\_shm\_mut}     & \yes & \no  & \no  & \no  & \no  & \no  \\
  \fn{zd\_shm\_mut\_drop}            & \yes & \no  & \no  & \no  & \no  & \no  \\
  \midrule
  \textbf{Totals}                    & \textbf{34} & \textbf{14} & \textbf{8} & \textbf{12} & \textbf{6} & \textbf{8} \\
\end{longtable}
}

\noindent
\textbf{Key:} P1 = \code{static inline} functions, P2 = \code{\_Generic}
macro dispatch, P3 = options structs, P4 = opaque \code{sizeof},
P5 = closure callbacks, P6 = loaning (const/mut enforcement).

\smallskip\noindent
\textbf{Note:} Loan functions are no longer marked P1. The per-type loan
functions (\fn{z\_config\_loan}, \fn{z\_bytes\_loan}, etc.) are exported
symbols in \file{libzenohc.so}---they are not \code{static inline}. P6
marks the 8 explicit loan wrappers whose primary purpose is const/mut
enforcement and macro bypass for the borrowing protocol.


% ══════════════════════════════════════════════════════════════════════════════
\clearpage
\part{Conformance Summary}
\label{part:conformance}
% ══════════════════════════════════════════════════════════════════════════════

\section{C Best Practices}

\begin{table}[ht]
\centering\small
\begin{tabularx}{\textwidth}{lcL}
  \toprule
  \textbf{Practice} & \textbf{Status} & \textbf{Notes} \\
  \midrule
  All exports use \code{FFI\_PLUGIN\_EXPORT}      & \PASS    & \code{visibility("default")} \\
  \code{C\_VISIBILITY\_PRESET hidden}             & \PASS    & Only marked symbols visible \\
  Consistent \code{zd\_} namespace prefix         & \PASS    & No collisions with \code{z\_}/\code{zc\_} \\
  Return code checking                             & \PARTIAL & F4: \fn{z\_encoding\_from\_str} \\
  Memory cleanup on all paths                      & \PASS    & Closures dropped on failure \\
  No undefined behavior                            & \PASS    & All pointer arithmetic bounded \\
  Const correctness                                & \PASS    & \code{const} on loan params \\
  Feature guards for optional functionality        & \PASS    & SHM behind \code{\#ifdef} \\
  \bottomrule
\end{tabularx}
\end{table}

\section{Dart FFI Best Practices}

\begin{table}[ht]
\centering\small
\begin{tabularx}{\textwidth}{lcL}
  \toprule
  \textbf{Practice} & \textbf{Status} & \textbf{Notes} \\
  \midrule
  Lazy singleton for bindings             & \PASS & Initialized on first access \\
  \fn{Dart\_InitializeApiDL} before NativePort & \PASS & Called in \fn{\_initBindings} \\
  \code{ReceivePort} for async callbacks  & \PASS & Subscriber, matching, scout \\
  \code{calloc.allocate(sizeof)} pattern  & \PASS & Consistent across all types \\
  Ownership tracking              & \PASS & All wrapper classes \\
  \code{try/finally} for temp allocations & \PASS & \fn{\_withKeyExpr}, encoding \\
  Idempotent \code{close()}/\code{dispose()} & \PASS & All classes \\
  Explicit length for strings & \PASS & \code{toDartString(length:)} \\
  \bottomrule
\end{tabularx}
\end{table}

\section{Architecture Conformance}

\begin{table}[ht]
\centering\small
\begin{tabularx}{\textwidth}{lcL}
  \toprule
  \textbf{Principle} & \textbf{Status} & \textbf{Notes} \\
  \midrule
  Shim is mechanical (no business logic) & \PASS & 1--5 zenoh-c calls per function \\
  Single-load library pattern            & \PASS & OS resolves \code{libzenohc.so} \\
  Options structs internalized           & \PASS & Dart never touches \code{z\_*\_options\_t} \\
  Move semantics handled in C            & \PASS & Dart never calls \code{z\_*\_move} \\
  Callback bridge via NativePort         & \PASS & Thread-safe, copies data \\
  Sentinel values for optional params    & \PASS & NULL for strings, $-1$ for enums \\
  \bottomrule
\end{tabularx}
\end{table}


% ══════════════════════════════════════════════════════════════════════════════
\clearpage
\part{Pattern Summary}
\label{part:summary}
% ══════════════════════════════════════════════════════════════════════════════

\begin{table}[ht]
\centering\footnotesize
\begin{tabularx}{\textwidth}{c >{\raggedright\arraybackslash}p{2.2cm} LL}
  \toprule
  \textbf{Pat.} & \textbf{What It Is} & \textbf{Why Dart Cannot Call It} & \textbf{How the Shim Resolves It} \\
  \midrule
  \textbf{1} & 56 \code{static inline} move functions &
    \code{static} = internal linkage; \code{inline} = no symbol.
    \code{nm~-D} confirms absence. ffigen permanently excludes them. &
    Shim calls them at compile time. C compiler inlines the cast into the
    exported function. \\[6pt]
  \textbf{2} & 4 \code{\_Generic} macros (25--56 branches) &
    Preprocessor construct. No symbol. Dispatches to type-specific
    functions at compile time. &
    Shim calls monomorphic functions directly. \\[6pt]
  \textbf{3} & Options structs with fields requiring~P1 &
    \code{sizeof} unknown. Fields require \fn{z\_*\_move}. Layout is
    version-specific. &
    Shim allocates, initializes, sets fields, and passes---all in one
    exported function. \\[6pt]
  \textbf{4} & Opaque type sizes (8--2008~bytes) &
    Dart FFI has no \code{sizeof()} for foreign types. Types mapped to
    \code{Opaque}. &
    Shim exports \fn{zd\_*\_sizeof()} functions. Dart queries at runtime. \\[6pt]
  \textbf{5} & Closure callbacks on Tokio threads &
    \fn{Pointer.from\-Function} cannot capture state. Non-Dart thread~= UB. &
    NativePort bridge: C callback posts \code{Dart\_CObject} via
    \fn{Dart\_Post\-CObject\_DL}. \\[6pt]
  \textbf{6} & Loaning protocol (const/mut borrowing) &
    Per-type functions \textbf{are} exported (\code{nm~-D} confirms), but
    \code{z\_loan}/\code{z\_loan\_mut} macros are not~(P2).
    Dart's \code{Pointer<Opaque>} erases const qualifier. &
    Shim wraps each loan with explicit naming
    (\fn{zd\_*\_loan} vs \fn{zd\_*\_loan\_mut}).
    Const enforced at C compiler level. \\
  \bottomrule
\end{tabularx}
\caption{Six FFI barrier patterns and their C shim resolutions}
\label{tab:pattern-summary}
\end{table}

\begin{patternbox}
The six patterns are not independent barriers --- they \textbf{compound}.
A single \fn{z\_put} call requires Pattern~4 (to allocate the payload),
Pattern~6 (to loan the session and key expression), Pattern~1 (to move the
payload), and Pattern~3 (to create and pass options). A subscriber
declaration adds Pattern~5 (callback bridge). The C shim collapses all of
these into single FFI-callable functions with simple scalar parameters.
\end{patternbox}


% ══════════════════════════════════════════════════════════════════════════════
\clearpage
\appendix
\part*{Appendices}
\addcontentsline{toc}{part}{Appendices}
% ══════════════════════════════════════════════════════════════════════════════

\section{Findings Summary}
\label{app:findings}

\begin{table}[ht]
\centering\small
\begin{tabularx}{\textwidth}{clL}
  \toprule
  \textbf{ID} & \textbf{Severity} & \textbf{Description} \\
  \midrule
  F1 & Low    & \fn{z\_close} return value ignored in \fn{zd\_close\_session} \\
  F2 & Low    & Payload passes through \fn{z\_bytes\_to\_string} in subscriber callback \\
  F3 & Low    & \code{ZBytes.fromUint8List} element-by-element copy \\
  F4 & Low    & \fn{z\_encoding\_from\_str} return code unchecked \\
  F5 & Low    & \fn{Zenoh.initLog} frees \code{malloc}'d string with \code{calloc.free} \\
  F6 & Medium & \fn{Zenoh.scout} ignores \fn{zd\_scout} return value \\
  F7 & Low    & \code{ZBytes.dispose()} leaks calloc wrapper when consumed \\
  \bottomrule
\end{tabularx}
\caption{All findings}
\label{tab:findings}
\end{table}


\section{References}
\label{app:references}

\subsection{Dart FFI Official Documentation}

All \code{api.dart.dev/stable/latest/} URLs resolve to the current stable SDK
(Dart 3.11.1 at time of writing). Using \code{/latest/} ensures links remain
valid as new Dart versions are released.

\begin{table}[ht]
\centering\small
\begin{tabularx}{\textwidth}{lL}
  \toprule
  \textbf{Document} & \textbf{URL} \\
  \midrule
  dart:ffi overview         & \url{https://api.dart.dev/stable/latest/dart-ffi/dart-ffi-library.html} \\
  C interop guide           & \url{https://dart.dev/interop/c-interop} \\
  DynamicLibrary class      & \url{https://api.dart.dev/stable/latest/dart-ffi/DynamicLibrary-class.html} \\
  DynamicLibrary.lookup     & \url{https://api.dart.dev/stable/latest/dart-ffi/DynamicLibrary/lookup.html} \\
  Pointer.fromFunction      & \url{https://api.dart.dev/stable/latest/dart-ffi/Pointer/fromFunction.html} \\
  Opaque class              & \url{https://api.dart.dev/stable/latest/dart-ffi/Opaque-class.html} \\
  Struct class              & \url{https://api.dart.dev/stable/latest/dart-ffi/Struct-class.html} \\
  sizeOf function           & \url{https://api.dart.dev/stable/latest/dart-ffi/sizeOf.html} \\
  NativeApi.initializeApiDLData & \url{https://api.dart.dev/stable/latest/dart-ffi/NativeApi/initializeApiDLData.html} \\
  ffigen package            & \url{https://pub.dev/packages/ffigen} \\
  \bottomrule
\end{tabularx}
\end{table}

\subsection{Bug Tracker Citations}

\begin{table}[ht]
\centering\small
\begin{tabularx}{\textwidth}{llL}
  \toprule
  \textbf{Issue} & \textbf{Status} & \textbf{Key Finding} \\
  \midrule
  \href{https://github.com/dart-lang/ffigen/issues/146}{ffigen\#146}
    & Closed & ffigen now \textbf{skips} inline functions \\
  \href{https://github.com/dart-lang/native/issues/459}{native\#459}
    & Partial & \code{static inline} permanently excluded \\
  \href{https://github.com/dart-lang/sdk/issues/48865}{sdk\#48865}
    & Open & Non-Dart thread callback = undefined behavior \\
  \href{https://github.com/dart-lang/sdk/issues/52689}{sdk\#52689}
    & Open & \code{NativeCallable.listener} limitations \\
  \bottomrule
\end{tabularx}
\end{table}

\subsection{zenoh-c Source Files}

\begin{table}[ht]
\centering\small
\begin{tabularx}{\textwidth}{lLr}
  \toprule
  \textbf{File} & \textbf{Content} & \textbf{Lines} \\
  \midrule
  \file{zenoh\_macros.h}   & \code{static inline} moves, \code{\_Generic}, C++ overloads & 1467 \\
  \file{zenoh\_commons.h}  & Exported function signatures, struct typedefs & \raisebox{0.5ex}{\texttildelow}5000 \\
  \file{zenoh\_opaque.h}   & Opaque type definitions with compile-time sizes & \raisebox{0.5ex}{\texttildelow}800 \\
  \file{zenoh\_dart.h}     & C shim header --- 62 exported functions & 543 \\
  \file{zenoh\_dart.c}     & C shim implementation & \raisebox{0.5ex}{\texttildelow}700 \\
  \bottomrule
\end{tabularx}
\end{table}

\subsection{ISO C Standards}

\begin{table}[ht]
\centering\small
\begin{tabular}{lll}
  \toprule
  \textbf{Standard} & \textbf{Section} & \textbf{Content} \\
  \midrule
  ISO C11 & \S6.5.1.1 & \code{\_Generic} selection --- compile-time type dispatch \\
  ISO C11 & \S6.7.4   & \code{inline} function specifiers \\
  ISO C11 & \S6.2.2   & \code{static} linkage --- internal to translation unit \\
  \bottomrule
\end{tabular}
\end{table}

\subsection{Symbol Table Verification}

Commands run against \file{extern/zenoh-c/target/release/libzenohc.so}
(zenoh-c v1.7.2):

\begin{lstlisting}[style=shellstyle]
# Verify static inline functions are NOT exported
nm -D libzenohc.so | grep z_bytes_move     # (no output)
nm -D libzenohc.so | grep z_config_move    # (no output)

# Verify _Generic macros are NOT exported
nm -D libzenohc.so | grep "z_loan$"        # (no output)
nm -D libzenohc.so | grep "z_drop$"        # (no output)

# Verify real functions ARE exported
nm -D libzenohc.so | grep " T z_put$"
  000000000039e3f0 T z_put
nm -D libzenohc.so | grep " T z_open$"
  000000000039acb0 T z_open
nm -D libzenohc.so | grep " T z_config_drop$"
  000000000038f790 T z_config_drop
\end{lstlisting}


\end{document}
